\documentclass[11pt,a4paper]{article}
\usepackage[T1]{fontenc}
\usepackage[utf8]{inputenc}
\usepackage{lmodern,microtype}
\usepackage{geometry}
\geometry{top=25mm,bottom=27mm,left=24mm,right=24mm}
\usepackage{amsmath,amssymb,amsthm,mathtools,bm}
\usepackage{booktabs,array,longtable,enumitem,multicol,graphicx,float}
\usepackage[most]{tcolorbox}
\usepackage{fancyhdr,lastpage,xcolor,hyperref,cleveref,titlesec}
\definecolor{navy}{HTML}{142C47}
\definecolor{teal}{HTML}{147E8B}
\definecolor{ink}{HTML}{243447}
\definecolor{pale}{HTML}{EDF4F7}
\definecolor{line}{HTML}{CBD9E2}
\hypersetup{colorlinks=true,linkcolor=navy,citecolor=teal,urlcolor=teal,pdftitle={Quantum Mechanics I: Assignment II - Complete Mathematical Solutions},pdfauthor={Mathematical Physics Study Manuscript},pdfsubject={Finite potential wells, scattering, operator uncertainty and quantum measurements}}
\setlength{\parindent}{0pt}
\setlength{\parskip}{5pt}
\emergencystretch=2em
\setlength{\headheight}{15pt}
\pagestyle{fancy}
\fancyhf{}
\fancyhead[L]{\small\textsf{QUANTUM MECHANICS I \textbar{} ASSIGNMENT II}}
\fancyhead[R]{\small\textsf{M.Sc. PHYSICS}}
\fancyfoot[L]{\small\textsf{Complete mathematical solutions \textbullet{} teaching preprint}}
\fancyfoot[R]{\small\thepage\ / \pageref{LastPage}}
\renewcommand{\headrulewidth}{0.3pt}
\titleformat{\section}{\Large\bfseries\color{navy}}{\thesection.}{0.7em}{}
\titleformat{\subsection}{\large\bfseries\color{navy}}{\thesubsection}{0.7em}{}
\titleformat{\subsubsection}{\normalsize\bfseries\color{teal}}{}{0pt}{}
\titlespacing*{\section}{0pt}{18pt}{9pt}
\titlespacing*{\subsection}{0pt}{14pt}{6pt}
\tcbset{colback=pale,colframe=line,arc=1.2mm,boxrule=0.6pt,left=8pt,right=8pt,top=7pt,bottom=7pt}
\newtcolorbox{questionbox}{breakable,title={\footnotesize\sffamily ASSIGNMENT QUESTION},fonttitle=\bfseries,colback=white,colframe=teal,coltitle=teal,boxed title style={colback=white,colframe=white},attach boxed title to top left={yshift=-2mm,xshift=5mm}}
\newtcolorbox{resultbox}{breakable,colback=pale,colframe=teal,boxrule=.7pt}
\newcommand{\key}[1]{\begin{resultbox}\centering\(\displaystyle #1\)\end{resultbox}}
\newcommand{\dd}{\mathrm{d}}
\newcommand{\iu}{\mathrm{i}}
\newcommand{\ee}{\mathrm{e}}
\newcommand{\Hil}{\mathcal H}
\newcommand{\R}{\mathbb R}
\newcommand{\id}{\mathbb I}
\newcommand{\avg}[1]{\langle #1\rangle}
\newcommand{\Var}{\operatorname{Var}}
\newcommand{\Cov}{\operatorname{Cov}}
\newcommand{\Sch}{Schr\"odinger}
\newcommand{\bra}[1]{\langle #1|}
\newcommand{\ket}[1]{|#1\rangle}
\newcommand{\braket}[2]{\langle #1|#2\rangle}
\newcommand{\abs}[1]{\left|#1\right|}
\allowdisplaybreaks[2]
\begin{document}
\begin{titlepage}
\thispagestyle{empty}
\vspace*{17mm}
{\color{teal}\rule{\textwidth}{1.5pt}}\\[10mm]
{\sffamily\fontsize{29}{34}\selectfont\bfseries Quantum Mechanics--I:\\[2mm] Assignment II}\par
\vspace{6mm}
{\sffamily\Large Complete Worked Solutions from First Principles}\par
\vspace{11mm}
{\large\textbf{12 questions \textbullet{} Part A (5) and Part B (7)}}\par
\vspace{8mm}
\begin{tcolorbox}[colback=pale,colframe=teal]
\textbf{Course basis:} Department of Physics \& Nanotechnology, SRM Institute of Science and Technology, Kattankulathur. M.Sc. Physics / M.Sc. Applied Physics, Semester I. Quantum Mechanics--I (PPY25103J) / Quantum Mechanics (PAP25103J), Assignment--II.\\[4pt]
\textbf{Main topics:} Cauchy--Schwarz inequality, observable compatibility, infinite-well probabilities and superpositions, sequential measurement, generalized uncertainty relations, potential steps, tunneling, and the finite square well.
\end{tcolorbox}
\vspace{9mm}
{\Large\bfseries\color{navy} Mathematical Physics Study Manuscript}\par
\vspace{4mm}
{\normalsize Detailed symbolic derivations \textbullet{} boundary matching \textbullet{} verified numerical figures}\par
\vfill
{\large November 2026}\par
\vspace{3mm}
{\small arXiv-style \LaTeX{} preprint layout; not an arXiv submission or an original research claim.}\par
\vspace{8mm}
{\color{teal}\rule{\textwidth}{1.5pt}}
\end{titlepage}
\begin{abstract}
This pedagogical manuscript supplies step-by-step solutions to the five short-answer and seven long-answer questions of Quantum Mechanics--I Assignment II. Part A develops a direct verification of Cauchy--Schwarz for Legendre polynomials, commutator tests for compatible observables, exact position probabilities in an infinite well, and the energy statistics of an equal superposition. Its final reference, ``Griffiths Problem 3.27,'' omits an edition; we treat the verified second-edition sequential-measurement problem as a fully worked reference case and present separately verified first- and third-edition alternatives in an appendix. Part B proves the full \Sch--Robertson inequality and derives scattering amplitudes and flux coefficients for both above-threshold and below-threshold potential steps, the exact rectangular-barrier tunneling coefficient, the infinite-well spectrum, its position--momentum uncertainties, and the bound states of a finite symmetric well. The numerical finite-well example is computed at the dimensionless strength $P=4$, with explicitly normalized eigenfunctions and original plots of the first three probability densities. Operator-domain qualifications and normalization conventions are stated where they materially affect the mathematics. This is a study preprint, not evidence of new physical research or journal review.
\end{abstract}
\begingroup
\small
\setlength{\parskip}{0pt}
\tableofcontents
\endgroup
\clearpage
\section{Scope, conventions, and source fidelity}
The wording and ordering of the 12 questions follow the uploaded course file \emph{QM-AS-1-5.pdf}, specifically its Assignment--II portion. Because the supplied sheet states only ``Problem 3.27 Griffiths'' and does not identify the edition, the result for Part A, Problem 5 is explicitly conditional on the reference edition. The complete text of the textbook problems is not reproduced; concise descriptions and original derivations are used instead. Published variants are independently cross-checked against Griffiths editions \cite{griffiths1,griffiths2,griffiths3}.

We adopt $\hbar=h/(2\pi)$, particle mass $m>0$, $\iu^2=-1$, and the one-dimensional position-space scalar product
\begin{equation}
\langle\phi,\psi\rangle=\int \phi^*(x)\psi(x)\,\dd x,\qquad \|\psi\|^2=1.
\end{equation}
Where the particle moves freely in a region of constant potential, $\hat p=-\iu\hbar\,\dd/\dd x$ and $\hat H=-\hbar^2(\dd^2/\dd x^2)/(2m)+V$. The square-well width in Problems A3, A4, B5, B6 is $L$ on $0<x<L$. In B7 we write a \emph{symmetric finite well} on $-a<x<a$ (full width $2a$). The potential step at $x=0$ is semi-infinite in B2 and B3; B4 instead has a finite barrier on $0<x<a$.

For products of unbounded operators, commutator calculations are initially understood on a common suitable dense domain, e.g. smooth compactly supported functions, with the self-adjoint extensions specified when necessary. Pure plane waves are generalized scattering states; their reflection and transmission \emph{probabilities} are obtained as ratios of currents rather than by global $L^2$ normalization.

\section{Part A --- short problems (1--5)}
\subsection{A1. Verify Cauchy--Schwarz for two Legendre functions}
\begin{questionbox}
Verify the Schwarz inequality on $[-1,1]$ for $f_1(x)=x$ and $f_2(x)=\tfrac12(3x^2-1)$.
\end{questionbox}
\textbf{Step 1: Set up the inner product.} Use $\langle f,g\rangle=\int_{-1}^{1}f^*(x)g(x)\,\dd x$. The Cauchy--Schwarz inequality is
\begin{equation}\abs{\langle f_1,f_2\rangle}^2\leq\langle f_1,f_1\rangle\langle f_2,f_2\rangle.
\end{equation}
\textbf{Step 2: Compute each norm exactly.}
\begin{align}
\langle f_1,f_1\rangle&=\int_{-1}^{1}x^2\,\dd x=\left[\frac{x^3}{3}\right]_{-1}^{1}=\frac23,\\
\langle f_2,f_2\rangle&=\frac14\int_{-1}^{1}(9x^4-6x^2+1)\,\dd x\nonumber\\
&=\frac14\left(9\cdot\frac25-6\cdot\frac23+2\right)=\frac25.
\end{align}
\textbf{Step 3: Compute the mixed scalar product.}
\begin{equation}
\langle f_1,f_2\rangle=\frac12\int_{-1}^{1}(3x^3-x)\,\dd x=0,
\end{equation}
because the integrand is odd on a symmetric interval. Consequently
\begin{equation}
\abs{\langle f_1,f_2\rangle}^2=0\leq \frac23\cdot\frac25=\frac4{15}.
\end{equation}
\key{\boxed{0\leq 4/15\quad\text{(verified).}}}
\textbf{Interpretation.} $f_1=P_1$ and $f_2=P_2$ are distinct Legendre polynomials. Their orthogonality is a special case of $\int_{-1}^{1}P_\ell P_j\,\dd x=2\delta_{\ell j}/(2\ell+1)$. Equality in Cauchy--Schwarz would require linear dependence; these nonzero functions are linearly independent, so the inequality is strict.

\subsection{A2. Compatibility of momentum and Hamiltonian for a free particle}
\begin{questionbox}
Show that for a free particle in one dimension momentum and Hamiltonian are compatible observables, whereas position and Hamiltonian are not.
\end{questionbox}
\textbf{Step 1: Write the differential operators.} For a free particle on the full real line,
\begin{equation}
\hat p=-\iu\hbar\frac{\dd}{\dd x},\quad \hat x=x,\quad\hat H=\frac{\hat p^2}{2m}=-\frac{\hbar^2}{2m}\frac{\dd^2}{\dd x^2}.
\end{equation}
\textbf{Step 2: Calculate $[\hat p,\hat H]$ directly.} On a smooth test function $\phi$,
\begin{align}
\hat p\hat H\phi&=(-\iu\hbar)\left(-\frac{\hbar^2}{2m}\right)\phi'''=\frac{\iu\hbar^3}{2m}\phi''',\\
\hat H\hat p\phi&=-\frac{\hbar^2}{2m}(-\iu\hbar\phi')''=\frac{\iu\hbar^3}{2m}\phi'''.
\end{align}
Thus $[\hat p,\hat H]\phi=0$ and
\begin{equation}\boxed{[\hat p,\hat H]=0.}\end{equation}
As $\hat H=p^2/(2m)$ in the momentum representation, $\hat H$ and $\hat p$ are simultaneously diagonalized by the Fourier transform. Their generalized common eigenfunctions $\ee^{\iu px/\hbar}$ have $E=p^2/(2m)$; the energy $E>0$ corresponds to $p=\pm\sqrt{2mE}$.

\textbf{Step 3: Calculate $[\hat x,\hat H]$.} The product rule gives
\begin{align}
\hat x\hat H\phi&=-\frac{\hbar^2}{2m}x\phi'',\\
\hat H\hat x\phi&=-\frac{\hbar^2}{2m}(x\phi)''=-\frac{\hbar^2}{2m}(2\phi'+x\phi'').
\end{align}
Subtract:
\begin{equation}
[\hat x,\hat H]\phi=\frac{\hbar^2}{m}\phi'=\frac{\iu\hbar}{m}\hat p\phi.
\end{equation}
Therefore,
\key{\boxed{[\hat p,\hat H]=0,\qquad[\hat x,\hat H]=\frac{\iu\hbar}{m}\hat p\neq0.}}
\textbf{Step 4: Physical meaning.} Momentum and energy are compatible observables for free motion, while position and energy generally are not. The Heisenberg equations reinforce this result:
\begin{equation}
\frac{\dd\hat p}{\dd t}=\frac{1}{\iu\hbar}[\hat p,\hat H]=0,\qquad
\frac{\dd\hat x}{\dd t}=\frac{1}{\iu\hbar}[\hat x,\hat H]=\frac{\hat p}{m}.
\end{equation}
The first is momentum conservation; the second is the quantum analog of uniform rectilinear motion.

\subsection{A3. Finding the particle in the central half of an infinite box}
\begin{questionbox}
Find the probability of detecting a particle between $x=L/4$ and $x=3L/4$ in the ground state and the first excited state of a one-dimensional infinite square well.
\end{questionbox}
\textbf{Step 1: Derive the stationary eigenfunctions.} For $V=0$ in $(0,L)$, the \Sch{} equation reads $\psi''+k^2\psi=0$. The general solution $A\sin kx+B\cos kx$ obeys $\psi(0)=0$, forcing $B=0$, and $\psi(L)=0$, forcing $kL=n\pi$ for $n\geq1$. Normalization follows from
\begin{equation}
1=\abs{A}^2\int_0^L\sin^2(n\pi x/L)\,\dd x=\abs A^2\frac L2.
\end{equation}
Therefore
\begin{equation}
\psi_n(x)=\sqrt{\frac2L}\sin\frac{n\pi x}{L},\quad 0<x<L;\qquad E_n=\frac{n^2\pi^2\hbar^2}{2mL^2}.
\end{equation}
\textbf{Step 2: Apply Born's probability rule.} The probability of the specified interval is
\begin{equation}
P_n=\int_{L/4}^{3L/4}\abs{\psi_n(x)}^2\,\dd x=\frac2L\int_{L/4}^{3L/4}\sin^2\frac{n\pi x}{L}\,\dd x.
\end{equation}
Using $2\sin^2\theta=1-\cos 2\theta$ gives the antiderivative
\begin{equation}
P_n=\left[\frac{x}{L}-\frac{\sin(2n\pi x/L)}{2n\pi}\right]_{L/4}^{3L/4}
=\frac12-\frac{\sin(3n\pi/2)-\sin(n\pi/2)}{2n\pi}.
\end{equation}
\textbf{Step 3: Ground state $n=1$.} Since $\sin(3\pi/2)=-1$ and $\sin(\pi/2)=1$,
\begin{equation}
P_1=\frac12-\frac{-1-1}{2\pi}=\boxed{\frac12+\frac1\pi\simeq0.81831}.
\end{equation}
\textbf{Step 4: First excited state $n=2$.} Since $\sin(3\pi)=\sin\pi=0$,
\begin{equation}
P_2=\frac12-\frac{0-0}{4\pi}=\boxed{\frac12}.
\end{equation}
\key{\boxed{P_{\text{ground}}=81.83\%,\qquad P_{\text{first excited}}=50\%.}}
The ground-state density peaks at $L/2$, whereas the first excited state has a node there, explaining the different central-interval probabilities.

\subsection{A4. Equal superposition of two infinite-well eigenstates}
\begin{questionbox}
A particle in an infinite square well of length $L$ is in the state $\psi(x)=A[\psi_1(x)+\psi_2(x)]$. Determine $A$, $\langle H\rangle$ and the energy uncertainty $\sigma_H$.
\end{questionbox}
The normalized eigenstates satisfy $\langle\psi_m,\psi_n\rangle=\delta_{mn}$ and $\hat H\psi_n=E_n\psi_n$, with $E_n=n^2\pi^2\hbar^2/(2mL^2)$.
\textbf{Step 1: Normalization.}
\begin{align}
1&=\langle\psi,\psi\rangle
=\abs A^2\left(\langle\psi_1,\psi_1\rangle+\langle\psi_1,\psi_2\rangle+\langle\psi_2,\psi_1\rangle+\langle\psi_2,\psi_2\rangle\right)\\
&=2\abs A^2.
\end{align}
Thus $\abs A=1/\sqrt2$; choosing the arbitrary global phase to be zero, $A=1/\sqrt2$.
\textbf{Step 2: Expectation of energy.} Using $\hat H\psi_n=E_n\psi_n$,
\begin{align}
\langle H\rangle&=\frac12\langle\psi_1+\psi_2,\hat H(\psi_1+\psi_2)\rangle\\
&=\frac12\langle\psi_1+\psi_2,E_1\psi_1+E_2\psi_2\rangle=\frac{E_1+E_2}{2}.
\end{align}
Since $E_2=4E_1$,
\begin{equation}\boxed{\langle H\rangle=\frac52E_1=\frac{5\pi^2\hbar^2}{4mL^2}.}\end{equation}
\textbf{Step 3: Second moment and variance.} Acting twice with the Hamiltonian gives $\hat H^2\psi_n=E_n^2\psi_n$, hence
\begin{align}
\langle H^2\rangle&=\frac{E_1^2+E_2^2}{2}=\frac{17}{2}E_1^2=\frac{17\pi^4\hbar^4}{8m^2L^4},\\
\sigma_H^2&=\langle H^2\rangle-\langle H\rangle^2=\left(\frac{17}{2}-\frac{25}{4}\right)E_1^2=\frac94E_1^2.
\end{align}
The nonnegative standard deviation is $\sigma_H=(3/2)E_1$.
\key{\boxed{A=\frac1{\sqrt2},\qquad\langle H\rangle=\frac{5\pi^2\hbar^2}{4mL^2},\qquad\sigma_H=\frac{3\pi^2\hbar^2}{4mL^2}.}}
\textbf{Interpretation.} The energy probabilities are $P(E_1)=P(E_2)=1/2$, so $\sigma_H=(E_2-E_1)/2$. For the time-evolving superposition
\begin{equation}
\Psi(x,t)=\frac{\psi_1(x)\ee^{-\iu E_1t/\hbar}+\psi_2(x)\ee^{-\iu E_2t/\hbar}}{\sqrt2},
\end{equation}
the position probability density oscillates through interference, but the energy distribution and its standard deviation remain time independent.

\subsection{A5. Griffiths Problem 3.27 --- verified second-edition interpretation}
\begin{questionbox}
The assignment reads only ``Problem 3.27 Griffiths,'' without stating the edition. Here we solve the verified \emph{second-edition (2005)} version on sequential measurements. A full solution to the distinct first-edition and third-edition versions is supplied in \Cref{app:griffiths_variants}.
\end{questionbox}
\textbf{Edition note.} The second-edition printed problem is titled ``Sequential measurements'' and uses the two bases and coefficients below \cite{griffiths2}. The first edition instead asks for Gaussian-function Gram--Schmidt orthonormalization \cite{griffiths1}; the third edition asks for spectral decomposition of operators and operator exponentials \cite{griffiths3}. No edition can be established solely from the uploaded assignment.

Let the nondegenerate eigenbases of the observables $\hat A$ and $\hat B$ be $\{\ket{\psi_1},\ket{\psi_2}\}$ and $\{\ket{\phi_1},\ket{\phi_2}\}$, with
\begin{equation}
\ket{\psi_1}=\frac35\ket{\phi_1}+\frac45\ket{\phi_2},\qquad
\ket{\psi_2}=\frac45\ket{\phi_1}-\frac35\ket{\phi_2}.
\end{equation}
\textbf{Step 1: Postmeasurement state of $A$.} If the ideal projective measurement of $A$ produces $a_1$, then immediately afterward the state is $\boxed{\ket{\psi_1}}$.

\textbf{Step 2: Measurement of $B$.} From Born's rule and orthonormality of the $B$-basis,
\begin{align}
P(b_1\mid a_1)&=\abs{\braket{\phi_1}{\psi_1}}^2=\left(\frac35\right)^2=\frac9{25},\\
P(b_2\mid a_1)&=\abs{\braket{\phi_2}{\psi_1}}^2=\left(\frac45\right)^2=\frac{16}{25}.
\end{align}
\textbf{Step 3: Reverse the basis transformation.} In matrix form, $U=\frac15\bigl(\begin{smallmatrix}3&4\\4&-3\end{smallmatrix}\bigr)$; one checks $U^\dagger U=I$ and $U^{-1}=U$. Thus
\begin{equation}
\ket{\phi_1}=\frac35\ket{\psi_1}+\frac45\ket{\psi_2},\qquad
\ket{\phi_2}=\frac45\ket{\psi_1}-\frac35\ket{\psi_2}.
\end{equation}
If the unreported $B$-outcome was $b_1$, the final probability of $a_1$ is $9/25$; if it was $b_2$, that conditional probability is $16/25$. Because the $B$ measurement actually occurred, its unrecorded alternatives add as probabilities, not probability amplitudes:
\begin{align}
P(a_1\text{ on final }A\mid a_1\text{ initially})
&=\frac9{25}\cdot\frac9{25}+\frac{16}{25}\cdot\frac{16}{25}\\
&=\frac{81+256}{625}=\boxed{\frac{337}{625}=0.5392}.
\end{align}
\textbf{Density-operator verification.} The intermediate unread measurement produces
\begin{equation}
\rho_B=\frac9{25}\ket{\phi_1}\bra{\phi_1}+\frac{16}{25}\ket{\phi_2}\bra{\phi_2},
\quad P(a_1)=\langle\psi_1|\rho_B|\psi_1\rangle=\frac{337}{625}.
\end{equation}
If the measurement of $B$ had been omitted entirely, an immediate repeat measurement of $A$ would instead yield $a_1$ with probability one. The intermediate measurement disturbs the state.

\section{Part B --- extended problems (1--7)}
\subsection{B1. Generalized uncertainty relation and position--momentum application}
\begin{questionbox}
Derive the generalized uncertainty relation. Apply it to the position and momentum operators in one dimension.
\end{questionbox}
\textbf{Step 1: Define centered observables and variances.} Let $\hat A$ and $\hat B$ be self-adjoint observables in a normalized state $\ket\psi$. Define
\begin{equation}
\delta\hat A=\hat A-\langle A\rangle\id,\quad \delta\hat B=\hat B-\langle B\rangle\id,\qquad
(\Delta A)^2=\langle(\delta\hat A)^2\rangle,\quad(\Delta B)^2=\langle(\delta\hat B)^2\rangle.
\end{equation}
For observables with finite variances, these expressions are nonnegative because they are squared Hilbert-space norms.

\textbf{Step 2: Apply Cauchy--Schwarz to fluctuation vectors.} Let
\begin{equation}
\ket u=\delta\hat A\ket\psi,\qquad \ket v=\delta\hat B\ket\psi.
\end{equation}
The Cauchy--Schwarz inequality states $\langle u,u\rangle\langle v,v\rangle\geq\abs{\langle u,v\rangle}^2$. Therefore
\begin{equation}\label{eq:cs_unc}
\boxed{(\Delta A)^2(\Delta B)^2\geq\abs{\langle\delta\hat A\,\delta\hat B\rangle}^2.}
\end{equation}
\textbf{Step 3: Separate symmetric and antisymmetric operator products.} Define the commutator and anticommutator
\begin{equation}
[\hat A,\hat B]=\hat A\hat B-\hat B\hat A,\qquad\{\hat A,\hat B\}=\hat A\hat B+\hat B\hat A.
\end{equation}
Then
\begin{equation}
\delta\hat A\delta\hat B=\frac12\{\delta\hat A,\delta\hat B\}+\frac12[\delta\hat A,\delta\hat B]
=\frac12\{\delta\hat A,\delta\hat B\}+\frac12[\hat A,\hat B].
\end{equation}
Self-adjointness of the centered observables implies that $\{\delta\hat A,\delta\hat B\}$ is Hermitian (real expectation) and $[\hat A,\hat B]$ is anti-Hermitian (purely imaginary expectation). Thus the two pieces are orthogonal as real and imaginary components of a complex number:
\begin{align}
\abs{\langle\delta\hat A\delta\hat B\rangle}^2
=\frac14\abs{\langle\{\delta\hat A,\delta\hat B\}\rangle}^2
+\frac14\abs{\langle[\hat A,\hat B]\rangle}^2.
\end{align}
\textbf{Step 4: Obtain the Schr\"odinger--Robertson inequality.} Substitution into \eqref{eq:cs_unc} gives
\begin{equation}\label{eq:SR}
\boxed{(\Delta A)^2(\Delta B)^2\geq\frac14\abs{\langle[\hat A,\hat B]\rangle}^2+
\frac14\abs{\langle\{\delta\hat A,\delta\hat B\}\rangle}^2.}
\end{equation}
The second term is the square of the symmetrized covariance, $\Cov(A,B)=\tfrac12\langle\{\delta\hat A,\delta\hat B\}\rangle$. Omitting this nonnegative term gives the more familiar Robertson relation
\begin{equation}\label{eq:robertson}
\boxed{\Delta A\,\Delta B\geq\frac12\abs{\langle[\hat A,\hat B]\rangle}.}
\end{equation}

\textbf{Step 5: Calculate the canonical commutator.} In one dimension,
\begin{equation}\hat x\phi=x\phi,\qquad\hat p\phi=-\iu\hbar\phi',\end{equation}
so on a smooth test function
\begin{align}
[\hat x,\hat p]\phi&=x(-\iu\hbar\phi')-[-\iu\hbar(x\phi)']\\
&=-\iu\hbar x\phi'+\iu\hbar(\phi+x\phi')=\iu\hbar\phi.
\end{align}
Hence $[\hat x,\hat p]=\iu\hbar\id$. Apply \eqref{eq:robertson}:
\begin{equation}
\Delta x\,\Delta p\geq\frac12\abs{\langle\iu\hbar\id\rangle}=\frac\hbar2.
\end{equation}
\key{\boxed{\Delta x\,\Delta p\geq\frac\hbar2,\qquad
(\Delta x)^2(\Delta p)^2\geq\frac{\hbar^2}4+\Cov(x,p)^2.}}
\textbf{Step 6: Exhibit a minimum-uncertainty state.} For $\sigma>0$, consider
\begin{equation}
\psi(x)=\frac{1}{(2\pi\sigma^2)^{1/4}}\exp\left[-\frac{(x-x_0)^2}{4\sigma^2}+\frac{\iu p_0 x}{\hbar}\right].
\end{equation}
Its Gaussian probability density has $\langle x\rangle=x_0$, $\Delta x=\sigma$, and direct differentiation yields
\begin{equation}
(\hat p-p_0)\psi=\frac{\iu\hbar(x-x_0)}{2\sigma^2}\psi,\quad
\langle p\rangle=p_0,\quad (\Delta p)^2=\frac{\hbar^2}{4\sigma^2}.
\end{equation}
Thus $\Delta x\Delta p=\hbar/2$: the lower bound is attainable for a suitable Gaussian state. Equality in the full Schr\"odinger--Robertson inequality occurs exactly when the two centered fluctuation vectors are linearly dependent; equality in the simpler Robertson form additionally requires zero symmetrized covariance.

\subsection{B2. Potential step: energy above the step ($E>V_0$)}
\begin{questionbox}
Derive the reflection and transmission probabilities for a particle with $E>V_0$ incident on a potential step of height $V_0$ at the origin.
\end{questionbox}
Let $V(x)=0$ for $x<0$ and $V(x)=V_0>0$ for $x\geq0$. With $E>V_0$, both regional kinetic energies are positive. Define
\begin{equation}
k_1=\frac{\sqrt{2mE}}{\hbar},\qquad k_2=\frac{\sqrt{2m(E-V_0)}}{\hbar}.
\end{equation}
\textbf{Step 1: Solve the two regional \Sch{} equations.} In Region I, $\psi_I''+k_1^2\psi_I=0$; in Region II, $\psi_{II}''+k_2^2\psi_{II}=0$. An incident wave from the left has the form
\begin{equation}
\psi_I=A\ee^{\iu k_1x}+B\ee^{-\iu k_1x},\qquad
\psi_{II}=C\ee^{\iu k_2x}.
\end{equation}
The coefficient of $\ee^{-\iu k_2x}$ is zero because there is no incoming particle from the right.

\textbf{Step 2: Match the wavefunction and derivative.} The potential is finite at $x=0$, so both $\psi$ and $\psi'$ are continuous:
\begin{equation}
A+B=C,\qquad \iu k_1(A-B)=\iu k_2C.
\end{equation}
Eliminate $C$ using $C=A+B$:
\begin{equation}
k_1(A-B)=k_2(A+B)\quad\Longrightarrow\quad(k_1-k_2)A=(k_1+k_2)B.
\end{equation}
Thus the reflection and transmission \emph{amplitude ratios} are
\begin{equation}
\boxed{r=\frac BA=\frac{k_1-k_2}{k_1+k_2},\qquad t=\frac CA=1+r=\frac{2k_1}{k_1+k_2}.}
\end{equation}
\textbf{Step 3: Derive current coefficients.} From the probability continuity equation,
\begin{equation}\label{eq:current}
j=\frac{\hbar}{2m\iu}\left(\psi^*\psi'-\psi(\psi^*)'\right).
\end{equation}
For a right-moving plane wave $F\ee^{\iu kx}$, $j=(\hbar k/m)\abs F^2$; for a left-moving plane wave, $j=-(\hbar k/m)\abs F^2$. Thus
\begin{equation}
j_{\rm inc}=\frac{\hbar k_1}{m}\abs A^2,\quad
j_{\rm ref}=-\frac{\hbar k_1}{m}\abs B^2,\quad
j_{\rm tr}=\frac{\hbar k_2}{m}\abs C^2.
\end{equation}
The \emph{probabilities} are ratios of the magnitudes of outgoing currents to incoming current:
\begin{align}
R&=\frac{\abs{j_{\rm ref}}}{j_{\rm inc}}=\abs{\frac BA}^2=\boxed{\frac{(k_1-k_2)^2}{(k_1+k_2)^2}},\\
T&=\frac{j_{\rm tr}}{j_{\rm inc}}=\frac{k_2}{k_1}\abs{\frac CA}^2=\boxed{\frac{4k_1k_2}{(k_1+k_2)^2}}.
\end{align}
\textbf{Step 4: Write the results in energy variables.} The common factor $\sqrt{2m}/\hbar$ cancels:
\begin{equation}
\boxed{R=\left(\frac{\sqrt E-\sqrt{E-V_0}}{\sqrt E+\sqrt{E-V_0}}\right)^2,\quad
T=\frac{4\sqrt{E(E-V_0)}}{(\sqrt E+\sqrt{E-V_0})^2}.}
\end{equation}
\textbf{Step 5: Verify flux conservation.}
\begin{equation}
R+T=\frac{(k_1-k_2)^2+4k_1k_2}{(k_1+k_2)^2}=\frac{(k_1+k_2)^2}{(k_1+k_2)^2}=1.
\end{equation}
\key{\boxed{R+T=1.}}
\textbf{Interpretation and limits.} Quantum reflection persists for $E>V_0$ if $k_1\ne k_2$, in contrast to a classical particle. If $V_0\to0$, $R\to0,T\to1$; if $E\to V_0^+$, $R\to1,T\to0$. In the high-energy limit, $R\simeq V_0^2/(16E^2)$. Crucially, $T\ne\abs t^2$ when incident and transmitted velocities differ.

\subsection{B3. Semi-infinite potential step: energy below the step ($E<V_0$)}
\begin{questionbox}
Calculate the reflection and transmission probabilities for a particle with $E<V_0$ incident on a potential step of height $V_0$ at the origin.
\end{questionbox}
Here $V(x)=0$ on $x<0$ and $V(x)=V_0$ on $x\geq0$, with $0<E<V_0$. Define
\begin{equation}
k=\frac{\sqrt{2mE}}{\hbar},\qquad\kappa=\frac{\sqrt{2m(V_0-E)}}{\hbar}.
\end{equation}
\textbf{Step 1: Regional differential equations.} For $x<0$, $\psi_I''+k^2\psi_I=0$, hence $\psi_I=A\ee^{\iu kx}+B\ee^{-\iu kx}$. For $x>0$, $\psi_{II}''-\kappa^2\psi_{II}=0$, hence generally $C\ee^{-\kappa x}+D\ee^{\kappa x}$. Boundedness as $x\to+\infty$ requires $D=0$. Therefore
\begin{equation}
\boxed{\psi_I=A\ee^{\iu kx}+B\ee^{-\iu kx},\qquad\psi_{II}=C\ee^{-\kappa x}.}
\end{equation}
\textbf{Step 2: Match at the origin.}
\begin{equation}
A+B=C,\qquad \iu k(A-B)=-\kappa C.
\end{equation}
Substituting the first condition into the second yields
\begin{equation}
(\iu k+\kappa)A=(\iu k-\kappa)B,
\end{equation}
so
\begin{equation}
\boxed{r=\frac BA=\frac{k-\iu\kappa}{k+\iu\kappa},\qquad \frac CA=1+r=\frac{2k}{k+\iu\kappa}.}
\end{equation}
\textbf{Step 3: Reflection probability.} Using the same incident and reflected velocities,
\begin{equation}
R=\abs r^2=\frac{\abs{k-\iu\kappa}^2}{\abs{k+\iu\kappa}^2}=\frac{k^2+\kappa^2}{k^2+\kappa^2}=\boxed{1}.
\end{equation}
\textbf{Step 4: Transmission current.} For $\psi_{II}=C\ee^{-\kappa x}$, its derivative is $\psi_{II}'=-\kappa\psi_{II}$, with real $\kappa$. Equation \eqref{eq:current} gives
\begin{equation}
j_{II}=\frac{\hbar}{2m\iu}[-\kappa\abs{\psi_{II}}^2+\kappa\abs{\psi_{II}}^2]=0.
\end{equation}
Hence $T=j_{II}/j_{\rm inc}=0$, despite $C\ne0$.
\key{\boxed{R=1,\qquad T=0,\qquad R+T=1.}}
\textbf{Step 5: Penetration depth.} The amplitude decays as $\ee^{-\kappa x}$. Its $1/e$ amplitude penetration depth is
\begin{equation}
\boxed{\delta=\frac1\kappa=\frac\hbar{\sqrt{2m(V_0-E)}}};\qquad
\abs{\psi_{II}(x)}^2=\abs C^2\ee^{-2\kappa x}.
\end{equation}
The probability-density decay length is $1/(2\kappa)$. The amplitude ratio $\abs{C/A}^2=4E/V_0$ is not a transmission probability; it can exceed unity because it is not a flux ratio. The finite forbidden-region case with nonzero transmission is treated next.

\subsection{B4. Exact tunneling through a finite rectangular barrier}
\begin{questionbox}
Derive the tunneling probability for a particle with $0<E<V_0$ incident on a finite potential barrier of height $V_0$ and width $a$.
\end{questionbox}
The potential is $V(x)=0$ for $x<0$, $V_0$ for $0<x<a$, and $0$ for $x>a$. Define
\begin{equation}
k=\frac{\sqrt{2mE}}{\hbar},\qquad \kappa=\frac{\sqrt{2m(V_0-E)}}{\hbar},\qquad k,\kappa>0.
\end{equation}
\begin{figure}[htbp]\centering
\includegraphics[width=.94\linewidth]{figures/step_and_barrier.pdf}
\caption{A semi-infinite step (left) and finite rectangular barrier (right). Only the finite barrier supports an outgoing right-moving current in the far-right region when $E<V_0$.}
\label{fig:steps}
\end{figure}
\textbf{Step 1: Solve the regional equations.} The \Sch{} equation gives propagating solutions outside and exponential solutions inside:
\begin{equation}\label{eq:barrier_waves}
\boxed{\begin{aligned}
\psi_I&=A\ee^{\iu kx}+B\ee^{-\iu kx}&&(x<0),\\
\psi_{II}&=C\ee^{\kappa x}+D\ee^{-\kappa x}&&(0<x<a),\\
\psi_{III}&=F\ee^{\iu kx}&&(x>a).
\end{aligned}}
\end{equation}
The $\ee^{-\iu kx}$ term is absent in Region III because there is no incidence from the right. \emph{Both} exponentials must be retained in Region II: there are two interfaces to match.

\textbf{Step 2: Write the four boundary equations.} Matching $\psi$ and $\psi'$ at $x=0$ gives
\begin{equation}
A+B=C+D,\qquad \iu k(A-B)=\kappa(C-D).
\end{equation}
At $x=a$ this gives
\begin{equation}
C\ee^{\kappa a}+D\ee^{-\kappa a}=F\ee^{\iu ka},\quad
\kappa(C\ee^{\kappa a}-D\ee^{-\kappa a})=\iu kF\ee^{\iu ka}.
\end{equation}
\textbf{Step 3: Solve the interface at $a$.} Add and subtract the last equations to isolate $C\ee^{\kappa a}$ and $D\ee^{-\kappa a}$:
\begin{align}
C&=\frac F2\ee^{\iu ka-\kappa a}\left(1+\frac{\iu k}{\kappa}\right),\\
D&=\frac F2\ee^{\iu ka+\kappa a}\left(1-\frac{\iu k}{\kappa}\right).
\end{align}
This produces the hyperbolic representation inside the barrier,
\begin{equation}\label{eq:barrier_hyperbolic}
\psi_{II}(x)=F\ee^{\iu ka}\left[\cosh\kappa(x-a)+\frac{\iu k}{\kappa}\sinh\kappa(x-a)\right].
\end{equation}
\textbf{Step 4: Evaluate at $x=0$.} The function and its derivative there are
\begin{align}
A+B&=F\ee^{\iu ka}\left[\cosh(\kappa a)-\frac{\iu k}{\kappa}\sinh(\kappa a)\right],\\
\iu k(A-B)&=F\ee^{\iu ka}\left[-\kappa\sinh(\kappa a)+\iu k\cosh(\kappa a)\right].
\end{align}
Dividing the latter by $\iu k$ and adding to the former eliminates $B$:
\begin{align}
2A&=F\ee^{\iu ka}\left[2\cosh(\kappa a)+\iu\left(\frac\kappa k-\frac k\kappa\right)\sinh(\kappa a)\right].
\end{align}
Thus the exact transmission amplitude is
\begin{equation}\label{eq:tbarrier}
\boxed{t=\frac FA=\frac{\ee^{-\iu ka}}{\cosh(\kappa a)+\iu\dfrac{\kappa^2-k^2}{2k\kappa}\sinh(\kappa a)}}.
\end{equation}
Subtracting the boundary equations similarly gives the reflection amplitude
\begin{equation}
\boxed{r=\frac BA=\frac{-\iu\dfrac{k^2+\kappa^2}{2k\kappa}\sinh(\kappa a)}{\cosh(\kappa a)+\iu\dfrac{\kappa^2-k^2}{2k\kappa}\sinh(\kappa a)}}.
\end{equation}
\textbf{Step 5: Calculate flux ratios.} Because the potential is zero on both sides, incident and transmitted waves have the same wave number $k$:
\begin{equation}
j_{\rm inc}=\frac{\hbar k}{m}\abs A^2,\qquad j_{\rm tr}=\frac{\hbar k}{m}\abs F^2,\qquad T=\frac{j_{\rm tr}}{j_{\rm inc}}=\abs t^2.
\end{equation}
Taking the modulus squared of \eqref{eq:tbarrier}, noting $\abs{\ee^{-\iu ka}}^2=1$, yields
\begin{equation}
T=\left[\cosh^2(\kappa a)+\frac{(\kappa^2-k^2)^2}{4k^2\kappa^2}\sinh^2(\kappa a)\right]^{-1}.
\end{equation}
Use $\cosh^2 y=1+\sinh^2y$ and the identity
\begin{equation}
4k^2\kappa^2+(\kappa^2-k^2)^2=(k^2+\kappa^2)^2
\end{equation}
to get
\begin{equation}
T=\left[1+\frac{(k^2+\kappa^2)^2}{4k^2\kappa^2}\sinh^2(\kappa a)\right]^{-1}.
\end{equation}
Since $k^2+\kappa^2=2mV_0/\hbar^2$ and $4k^2\kappa^2=16m^2E(V_0-E)/\hbar^4$,
\begin{equation}\label{eq:Tbarrier}
\boxed{T=\left[1+\frac{V_0^2}{4E(V_0-E)}\sinh^2\left(\frac{a\sqrt{2m(V_0-E)}}\hbar\right)\right]^{-1}.}
\end{equation}
\textbf{Step 6: Obtain $R$ and verify conservation.} Let $Q=(k^2+\kappa^2)^2\sinh^2(\kappa a)/(4k^2\kappa^2)$. The reflection amplitude gives $R=Q/(1+Q)$ while $T=1/(1+Q)$. Thus $R+T=1$. Explicitly,
\begin{equation}
\boxed{R=\frac{V_0^2\sinh^2(\kappa a)}{4E(V_0-E)+V_0^2\sinh^2(\kappa a)},\qquad R+T=1.}
\end{equation}
Even inside the forbidden region, a superposition of exponentials can carry flux:
\begin{equation}
\psi_{II}=C\ee^{\kappa x}+D\ee^{-\kappa x}\implies
j_{II}=\frac{2\hbar\kappa}{m}\operatorname{Im}(CD^*)=\frac{\hbar k}{m}\abs F^2.
\end{equation}
\textbf{Step 7: Thick-barrier approximation.} For $\kappa a\gg1$, $\sinh^2(\kappa a)\simeq\ee^{2\kappa a}/4$; the exponential term dominates the denominator in \eqref{eq:Tbarrier}, so
\begin{equation}
\boxed{T\simeq\frac{16E(V_0-E)}{V_0^2}\ee^{-2\kappa a}
=\frac{16E(V_0-E)}{V_0^2}\exp\left(-\frac{2a\sqrt{2m(V_0-E)}}{\hbar}\right).}
\end{equation}
\textbf{Checks.} At $a\to0$, $T\to1$. At $a\to\infty$ with $E<V_0$, $T\to0$. At $E\to V_0^-$, use $\sinh(\kappa a)\sim\kappa a$ to obtain the finite limit $T\to[1+mV_0a^2/(2\hbar^2)]^{-1}$. The thick-barrier expansion is \emph{not} valid near this last limit. The exponent is consistent with the WKB action $2\hbar^{-1}\int_0^a\sqrt{2m(V_0-E)}\,\dd x=2\kappa a$.

\subsection{B5. Infinite square well: eigenenergies and normalized eigenfunctions}
\begin{questionbox}
Derive (i) the energy eigenvalues and (ii) the normalized eigenfunctions for a particle trapped in an infinite square well of length $L$.
\end{questionbox}
Let $V=0$ on $(0,L)$ and $V=+\infty$ outside. For an ideal infinite well the appropriate boundary conditions are $\psi(0)=\psi(L)=0$. Inside,
\begin{equation}\hat H=-\frac{\hbar^2}{2m}\frac{\dd^2}{\dd x^2},\qquad-\frac{\hbar^2}{2m}\psi''=E\psi.\end{equation}
\textbf{Step 1: Prove positivity of nontrivial eigenenergies.} Multiply by $\psi^*$, integrate over $(0,L)$, and integrate by parts:
\begin{align}
E\int_0^L\abs\psi^2\,\dd x
&=-\frac{\hbar^2}{2m}\int_0^L\psi^*\psi''\,\dd x\\
&=\frac{\hbar^2}{2m}\int_0^L\abs{\psi'}^2\,\dd x,
\end{align}
where the boundary term $[\psi^*\psi']_0^L$ is zero. For a nonzero function satisfying both Dirichlet conditions, the integral of $\abs{\psi'}^2$ is positive, hence $E>0$.

\textbf{Step 2: Solve the Sturm--Liouville equation.} Put $k=\sqrt{2mE}/\hbar$, so
\begin{equation}
\psi''+k^2\psi=0,\qquad\psi=A\sin kx+B\cos kx.
\end{equation}
The left condition gives $B=0$. The right condition gives $A\sin kL=0$. Since $A\ne0$, $\sin kL=0$, i.e. $kL=n\pi$. Nontrivial states require $n=1,2,\ldots$, so
\begin{equation}
\boxed{k_n=\frac{n\pi}{L},\qquad E_n=\frac{\hbar^2k_n^2}{2m}=\frac{n^2\pi^2\hbar^2}{2mL^2}=\frac{n^2h^2}{8mL^2}.}
\end{equation}
In wavelength language, the wavefunction fits $n$ half-wavelengths: $L=n\lambda_n/2$, or $\lambda_n=2L/n$. The energies are not equally spaced:
\begin{equation}
E_{n+1}-E_n=(2n+1)\frac{\pi^2\hbar^2}{2mL^2}.
\end{equation}
\textbf{Step 3: Normalize.} With $\psi_n=A_n\sin(n\pi x/L)$,
\begin{align}
1&=\abs{A_n}^2\int_0^L\sin^2\frac{n\pi x}{L}\,\dd x
=\frac{\abs{A_n}^2}{2}\left[x-\frac{L}{2n\pi}\sin\frac{2n\pi x}{L}\right]_0^L
=\frac{\abs{A_n}^2L}{2}.
\end{align}
Choosing $A_n>0$ yields $A_n=\sqrt{2/L}$, therefore
\begin{equation}\label{eq:infinite_eigen}
\boxed{\psi_n(x)=\begin{cases}\sqrt{2/L}\,\sin(n\pi x/L),&0<x<L,\\0,&\text{otherwise},\end{cases}\quad n=1,2,\ldots.}
\end{equation}
The full stationary state is $\Psi_n(x,t)=\psi_n(x)\ee^{-\iu E_nt/\hbar}$.

\textbf{Step 4: Check orthogonality explicitly.} For $m\ne n$, use $2\sin\alpha\sin\beta=\cos(\alpha-\beta)-\cos(\alpha+\beta)$:
\begin{align}
\langle\psi_m,\psi_n\rangle&=\frac2L\int_0^L\sin\frac{m\pi x}{L}\sin\frac{n\pi x}{L}\,\dd x\\
&=\frac1L\left[\frac{L\sin[(m-n)\pi x/L]}{(m-n)\pi}
-\frac{L\sin[(m+n)\pi x/L]}{(m+n)\pi}\right]_0^L=0.
\end{align}
With normalization for $m=n$ we obtain $\boxed{\langle\psi_m,\psi_n\rangle=\delta_{mn}}$.

\textbf{Step 5: Completeness and arbitrary states.} The Fourier sine theorem establishes that the normalized eigenfunctions form a complete orthonormal basis of $L^2(0,L)$. Thus, for any $f\in L^2(0,L)$,
\begin{equation}
f(x)=\sum_{n=1}^{\infty}c_n\psi_n(x)\quad\text{in }L^2,\qquad c_n=\sqrt{\frac2L}\int_0^L f(x)\sin\frac{n\pi x}{L}\,\dd x.
\end{equation}
The distributional resolution of the identity is
\begin{equation}
\sum_{n=1}^{\infty}\psi_n(x)\psi_n^*(x')=\delta(x-x'),\qquad 0<x,x'<L.
\end{equation}
An arbitrary initial normalized state evolves as $\Psi(x,t)=\sum_nc_n\psi_n(x)\ee^{-\iu E_nt/\hbar}$, with energy outcome probabilities $\abs{c_n}^2$.

\textbf{Step 6: Parity.} Because the well is centered at $L/2$, the relevant reflection maps $x\mapsto L-x$, not $x\mapsto-x$. By trigonometry,
\begin{equation}
\psi_n(L-x)=\sqrt{2/L}\sin[n\pi-n\pi x/L]=(-1)^{n+1}\psi_n(x).
\end{equation}
Hence odd $n$ states are even under reflection about $L/2$, even $n$ states are odd, and the $n$th wavefunction has $n-1$ interior nodes.

\textbf{Step 7: Zero-point energy and rigor.} The $n=0$ sine solution is identically zero. Directly at $E=0$, the solution $\psi=Ax+B$ with $\psi(0)=\psi(L)=0$ also forces $A=B=0$. Consequently
\begin{equation}
\boxed{E_1=\frac{\pi^2\hbar^2}{2mL^2}>0.}
\end{equation}
The sharp Poincar\'e inequality $\int_0^L\abs{\psi'}^2\,\dd x\geq(\pi^2/L^2)\int_0^L\abs\psi^2\,\dd x$ gives the variational bound $\langle H\rangle\geq E_1$ for any normalized $\psi\in H_0^1(0,L)$. The self-adjoint Dirichlet Hamiltonian has operator domain $H^2(0,L)\cap H_0^1(0,L)$; all its spectral eigenvalues are positive, discrete, and simple.

\subsection{B6. Verify $\Delta x\,\Delta p_x\geq\hbar/2$ in the $n$th infinite-well state}
\begin{questionbox}
Verify the position--momentum uncertainty relation for the $n$th eigenstate of a particle trapped in an infinite square well of length $L$.
\end{questionbox}
Take the normalized real eigenfunction $\psi_n=\sqrt{2/L}\sin(n\pi x/L)$ on $(0,L)$, with $n\geq1$. Compute
\begin{equation}
(\Delta x)^2=\langle x^2\rangle-\langle x\rangle^2,\qquad
(\Delta p)^2=\langle p^2\rangle-\langle p\rangle^2.
\end{equation}
\textbf{Step 1: Mean position.} Using $2\sin^2\theta=1-\cos2\theta$ and setting $\alpha=2n\pi/L$,
\begin{align}
\langle x\rangle&=\frac2L\int_0^Lx\sin^2\frac{n\pi x}{L}\,\dd x\\
&=\frac1L\int_0^Lx\bigl(1-\cos\alpha x\bigr)\,\dd x
=\frac1L\left[\frac{L^2}{2}-\left(\frac{x\sin\alpha x}{\alpha}+\frac{\cos\alpha x}{\alpha^2}\right)_0^L\right]\\
&=\boxed{\frac L2},
\end{align}
because $\sin\alpha L=0$ and $\cos\alpha L=1$.

\textbf{Step 2: Mean squared position.}
\begin{equation}
\langle x^2\rangle=\frac1L\left[\int_0^Lx^2\,\dd x-\int_0^Lx^2\cos\alpha x\,\dd x\right].
\end{equation}
The second integral is evaluated by two integrations by parts:
\begin{align}
\int x^2\cos\alpha x\,\dd x&=\frac{x^2\sin\alpha x}{\alpha}-\frac2\alpha\int x\sin\alpha x\,\dd x,\\
\int x\sin\alpha x\,\dd x&=-\frac{x\cos\alpha x}{\alpha}+\frac{\sin\alpha x}{\alpha^2}.
\end{align}
Evaluate at $0,L$ to obtain
\begin{equation}
\int_0^Lx^2\cos\alpha x\,\dd x=\frac{2L}{\alpha^2}=\frac{L^3}{2n^2\pi^2}.
\end{equation}
Since $\int_0^Lx^2\,\dd x=L^3/3$,
\begin{equation}
\boxed{\langle x^2\rangle=L^2\left(\frac13-\frac1{2n^2\pi^2}\right).}
\end{equation}
\textbf{Step 3: Position variance.} Subtract $\langle x\rangle^2=L^2/4$:
\begin{equation}
\boxed{(\Delta x)^2=L^2\left(\frac1{12}-\frac1{2n^2\pi^2}\right),\qquad
\Delta x=L\sqrt{\frac1{12}-\frac1{2n^2\pi^2}}.}
\end{equation}

\textbf{Step 4: Mean momentum.} Using $\hat p=-\iu\hbar\,\dd/\dd x$,
\begin{align}
\langle p\rangle&=-\iu\hbar\int_0^L\psi_n(x)\psi_n'(x)\,\dd x
=-\frac{\iu\hbar}{2}\left[\psi_n^2(x)\right]_0^L=\boxed{0}.
\end{align}
\textbf{Step 5: Second momentum moment.} The eigenfunctions obey
\begin{equation}
-\hbar^2\psi_n''(x)=\frac{n^2\pi^2\hbar^2}{L^2}\psi_n(x).
\end{equation}
The kinetic-energy quadratic form gives
\begin{align}
\langle p^2\rangle&=\hbar^2\int_0^L\abs{\psi_n'}^2\,\dd x
=-\hbar^2\int_0^L\psi_n^*\psi_n''\,\dd x\\
&=\frac{n^2\pi^2\hbar^2}{L^2}\int_0^L\abs{\psi_n}^2\,\dd x
=\boxed{\frac{n^2\pi^2\hbar^2}{L^2}}.
\end{align}
The boundary term in the integration by parts vanishes by Dirichlet conditions. Therefore
\begin{equation}
\boxed{\Delta p=\frac{n\pi\hbar}{L}.}
\end{equation}
\textbf{Step 6: Uncertainty product.} Multiplication gives
\begin{align}
\Delta x\,\Delta p
&=\frac{n\pi\hbar}{L}\,L\sqrt{\frac1{12}-\frac1{2n^2\pi^2}}\\
&=\hbar\sqrt{\frac{n^2\pi^2-6}{12}}
=\boxed{\frac\hbar2\sqrt{\frac{n^2\pi^2}{3}-2}}.
\end{align}
For $n\geq1$, $n^2\pi^2\geq\pi^2>9$, so $(n^2\pi^2/3)-2>1$. Hence
\begin{equation}\boxed{\Delta x\,\Delta p>\frac\hbar2\quad\text{for every }n\geq1.}\end{equation}
\textbf{Step 7: Ground state and semiclassical limit.} At $n=1$,
\begin{equation}
\Delta x_1\Delta p_1=\hbar\sqrt{\frac{\pi^2-6}{12}}\simeq0.56786\hbar>0.5\hbar.
\end{equation}
For $n\to\infty$ at fixed $L,\hbar$,
\begin{equation}
\Delta x\longrightarrow\frac{L}{\sqrt{12}},\quad \Delta p\sim\frac{n\pi\hbar}{L},\quad
\Delta x\Delta p\sim\frac{n\pi\hbar}{\sqrt{12}}.
\end{equation}
The first limit equals the standard deviation of a \emph{uniform} classical position distribution on $(0,L)$, consistent with the rapidly oscillating $\sin^2(n\pi x/L)$ approaching $1/2$ under coarse graining. The product grows rather than approaching the lower bound.

\textbf{Operator-domain note.} On $L^2(0,L)$ the Dirichlet differential momentum operator is symmetric but is not itself self-adjoint with that domain. The expression above for $\langle p^2\rangle$ is rigorously a positive quadratic form. Alternatively extend the wavefunction by zero to $\R$ and use the self-adjoint full-line momentum operator: the zero extension lies in $H^1(\R)$, so $\langle p\rangle$ and $\|p\psi\|^2$ are meaningful; it need not lie in the domain of the full-line $p^2$. The usual Heisenberg variance therefore remains well defined. A direct Cauchy--Schwarz verification uses $u=(x-L/2)\psi_n$ and $v=-\iu\hbar\psi_n'$; integrating by parts gives $\langle u,v\rangle=\iu\hbar/2$ and $\|u\|\|v\|\geq\hbar/2$.

\subsection{B7. Finite symmetric square well: exact bound-state equations and plots}
\begin{questionbox}
Solve the time-independent \Sch{} equation for a particle trapped in a finite square well of height $U$ and obtain the wavefunctions. Plot the first three eigenfunctions and their probability densities. Because the assignment specifies no width or energy convention, we choose the standard symmetric well of depth $U$ and half-width $a$, and explicitly state the equivalent positive-height convention.
\end{questionbox}
\textbf{Step 1: Choose an energy convention.} Let
\begin{equation}\label{eq:Vfinite}
\boxed{V(x)=\begin{cases}-U,&\abs x<a,\\0,&\abs x\geq a,\end{cases}\qquad U>0.}
\end{equation}
The full well width is $2a$. Bound states have energy $-U<E<0$. The convention $V=0$ for $\abs x<a$ and $V=U$ outside differs only by the additive constant $U$ in energy; eigenfunctions and matching equations do not change.

\textbf{Step 2: Solve in three regions.} Introduce
\begin{equation}
k=\frac{\sqrt{2m(E+U)}}\hbar>0,\quad \kappa=\frac{\sqrt{-2mE}}\hbar>0,\quad k^2+\kappa^2=\frac{2mU}{\hbar^2}.
\end{equation}
The exterior equations $\psi''-\kappa^2\psi=0$ have only exponentially decaying admissible tails. The interior equation $\psi''+k^2\psi=0$ has oscillatory solutions. The general square-integrable form is
\begin{equation}
\psi(x)=\begin{cases}L_0\ee^{\kappa(x+a)},&x<-a,\\B\cos(kx)+C\sin(kx),&\abs x\leq a,\\R_0\ee^{-\kappa(x-a)},&x>a.\end{cases}
\end{equation}
\textbf{Step 3: Classify states by parity.} Since $V(-x)=V(x)$ and one-dimensional bound eigenvalues are nondegenerate, the eigenfunctions can be chosen even or odd. For an even state, set $C=0$ and $L_0=R_0=A_+\cos(ka)$; for an odd state, set $B=0$, $L_0=-A_-\sin(ka)$ and $R_0=A_-\sin(ka)$. Hence
\begin{equation}\label{eq:finite_even}
\boxed{\psi_+(x)=A_+\begin{cases}\cos(ka)\ee^{\kappa(x+a)},&x<-a,\\\cos(kx),&\abs x\leq a,\\\cos(ka)\ee^{-\kappa(x-a)},&x>a,\end{cases}}
\end{equation}
\begin{equation}\label{eq:finite_odd}
\boxed{\psi_-(x)=A_-\begin{cases}-\sin(ka)\ee^{\kappa(x+a)},&x<-a,\\\sin(kx),&\abs x\leq a,\\\sin(ka)\ee^{-\kappa(x-a)},&x>a.\end{cases}}
\end{equation}
\textbf{Step 4: Match at $x=a$.} The boundary-value construction already enforces continuity of $\psi$; derivative continuity gives
\begin{align}
\text{even:}\quad -A_+k\sin(ka)&=-A_+\kappa\cos(ka)\ \Longrightarrow\ \boxed{k\tan(ka)=\kappa},\\
\text{odd:}\quad A_-k\cos(ka)&=-A_-\kappa\sin(ka)\ \Longrightarrow\ \boxed{-k\cot(ka)=\kappa}.
\end{align}
The conditions at $x=-a$ then follow automatically from parity. The energy eigenvalues cannot generally be written in elementary closed form: they are fixed by roots of these two transcendental equations.

\textbf{Step 5: Dimensionless quantization.} Define the dimensionless variables
\begin{equation}
\boxed{z=ka,\quad q=\kappa a,\quad P=\frac{a\sqrt{2mU}}\hbar.}
\end{equation}
They satisfy $z^2+q^2=P^2$, with $0<z<P$ and $q=\sqrt{P^2-z^2}$. The allowed roots obey
\begin{equation}\label{eq:evenodd_finite}
\boxed{\begin{aligned}z\tan z&=\sqrt{P^2-z^2}&&\text{(even)},\\
-z\cot z&=\sqrt{P^2-z^2}&&\text{(odd)}.
\end{aligned}}
\end{equation}
For a root $z_n$, the physical energy is
\begin{equation}\label{eq:Efinitex}
\boxed{E_n=-U+\frac{\hbar^2z_n^2}{2ma^2}=-\frac{\hbar^2q_n^2}{2ma^2}.}
\end{equation}
With the positive-barrier convention ($0$ inside, $U$ outside), the same state has $E_n'=E_n+U=\hbar^2z_n^2/(2ma^2)$.

\textbf{Step 6: Normalize the even states.} From \eqref{eq:finite_even}, using evenness of $\abs\psi^2$,
\begin{align}
1&=2\abs{A_+}^2\left[\int_0^a\cos^2(kx)\,\dd x+\cos^2(ka)\int_a^\infty\ee^{-2\kappa(x-a)}\,\dd x\right]\\
&=\abs{A_+}^2\left[a+\frac{\sin(2ka)}{2k}+\frac{\cos^2(ka)}{\kappa}\right].
\end{align}
The even quantization condition gives $\tan(ka)=\kappa/k$. Thus $\sin(ka)\cos(ka)=k\kappa/(k^2+\kappa^2)$ and $\cos^2(ka)=k^2/(k^2+\kappa^2)$, so
\begin{equation}
\frac{\sin(2ka)}{2k}+\frac{\cos^2(ka)}{\kappa}
=\frac{\kappa}{k^2+\kappa^2}+\frac{k^2}{\kappa(k^2+\kappa^2)}=\frac1\kappa.
\end{equation}
Therefore
\begin{equation}
\boxed{A_+=\frac1{\sqrt{a+1/\kappa}}.}
\end{equation}
\textbf{Step 7: Normalize the odd states.} From \eqref{eq:finite_odd},
\begin{equation}
1=\abs{A_-}^2\left[a-\frac{\sin(2ka)}{2k}+\frac{\sin^2(ka)}{\kappa}\right].
\end{equation}
The odd condition gives $\cot(ka)=-\kappa/k$, $\sin(ka)\cos(ka)=-k\kappa/(k^2+\kappa^2)$, and $\sin^2(ka)=k^2/(k^2+\kappa^2)$. Hence the bracket is again $a+1/\kappa$:
\begin{equation}
\boxed{A_-=\frac1{\sqrt{a+1/\kappa}}.}
\end{equation}
Both parities are normalized by the same formula when evaluated with the $\kappa$ of the individual state. For each $n$ let $N_n=(a+1/\kappa_n)^{-1/2}$; the \emph{fully normalized} wavefunctions are then exactly \eqref{eq:finite_even} and \eqref{eq:finite_odd} with $A_\pm=N_n$.

\textbf{Step 8: Count the bound states.} A useful single root equation follows by setting $z=P\sin\theta$, $q=P\cos\theta$ for $0<\theta<\pi/2$. Even roots satisfy $z+\theta=(2j+1)\pi/2$; odd roots satisfy $z+\theta=(j+1)\pi$. Ordering both families from the ground state with index $n=0$ yields
\begin{equation}\label{eq:unifiedf}
\boxed{z_n+\arcsin(z_n/P)=\frac{(n+1)\pi}{2},\quad 0<z_n<P,\quad n=0,1,\ldots.}
\end{equation}
The left side increases strictly on $(0,P)$ and tends to $P+\pi/2$ at the right endpoint. Thus the $n$th bound state exists if and only if $(n+1)\pi/2<P+\pi/2$, or $P>n\pi/2$. A threshold with $q=0$ is excluded because the exterior solution would not decay. The exact number of negative-energy bound states is
\begin{equation}
\boxed{N_{\rm bound}=\left\lceil\frac{2P}{\pi}\right\rceil\qquad(P>0).}
\end{equation}
Three bound states exist exactly when
\begin{equation}
\boxed{P>\pi\quad\Longleftrightarrow\quad U>\frac{\pi^2\hbar^2}{2ma^2}.}
\end{equation}
Any attractive finite square well in one dimension has at least one bound state.

\textbf{Step 9: Numerically determine and plot the first three states.} The assignment supplies no numerical well parameters, so we choose $P=4$ solely for illustration. This lies in $\pi<P<3\pi/2$ and supports exactly three bound states. Let $\varepsilon_0=\hbar^2/(2ma^2)$, so $U=16\varepsilon_0$. Solve the monotone equation \eqref{eq:unifiedf} for $n=0,1,2$. The results are listed in \Cref{tab:finite_roots}; the roots were independently checked by direct substitution into the parity equations.
\begin{table}[htbp]\centering
\caption{First three finite-well bound states for $P=4$, with zero energy outside the well. The normalization uses the dimensionless coordinate $u=x/a$.}
\label{tab:finite_roots}
\begin{tabular}{@{}llllrr@{}}\toprule
State & Parity & $z_n=k_na$ & $q_n=\kappa_na$ & $E_n/\varepsilon_0$ & $\sqrt a\,N_n$\\\midrule
$n=0$ & even & $1.252353$ & $3.798896$ & $-14.431611$ & $0.889730$\\
$n=1$ & odd  & $2.474577$ & $3.142685$ & $-9.876470$  & $0.870983$\\
$n=2$ & even & $3.595305$ & $1.753221$ & $-3.073783$  & $0.797991$\\\bottomrule
\end{tabular}
\end{table}

Define $u=x/a$ and $\varphi_n(u)=\sqrt a\,\psi_n(au)$. The normalized dimensionless amplitudes are $\sqrt a\,N_n=(1+1/q_n)^{-1/2}$, and $\int_\R\abs{\varphi_n(u)}^2\,\dd u=1$. The original vector figures \ref{fig:finite_levels}--\ref{fig:finite_densities} use the exact piecewise functions with these numerical roots and show both interior oscillations and exterior exponential tails.
\begin{figure}[htbp]\centering\includegraphics[width=.88\linewidth]{figures/finite_well_levels.pdf}
\caption{Finite symmetric well with $P=4$ and three bound-state energy eigenvalues in units of $\varepsilon_0=\hbar^2/(2ma^2)$.}
\label{fig:finite_levels}
\end{figure}
\begin{figure}[htbp]\centering\includegraphics[width=.88\linewidth]{figures/finite_well_wavefunctions.pdf}
\caption{First three normalized finite-well eigenfunctions: even ground state ($n=0$), odd first excited state ($n=1$), and even second excited state ($n=2$). Dashed lines mark $x=\pm a$.}
\label{fig:finite_wavefunctions}
\end{figure}
\begin{figure}[htbp]\centering\includegraphics[width=.88\linewidth]{figures/finite_well_densities.pdf}
\caption{Probability densities corresponding to \Cref{fig:finite_wavefunctions}. Each curve integrates to unity on the whole real line. The less tightly bound $n=2$ state has the largest exterior penetration probability.}
\label{fig:finite_densities}
\end{figure}

\textbf{Step 10: Find the exterior detection probability.} From the normalized even or odd tails, the exact probability outside the geometrical well is
\begin{align}
P_{\rm out}^{(+)}&=2N_n^2\cos^2(k_na)\int_a^\infty \ee^{-2\kappa_n(x-a)}\,\dd x=\frac{N_n^2\cos^2(k_na)}{\kappa_n},\\
P_{\rm out}^{(-)}&=\frac{N_n^2\sin^2(k_na)}{\kappa_n}.
\end{align}
Since $N_n^2=1/(a+1/\kappa_n)$, these can be written in dimensionless form as $\cos^2z_n/(q_n+1)$ (even) or $\sin^2z_n/(q_n+1)$ (odd). For the $P=4$ example,
\begin{equation}
\boxed{P_{\rm out,0}\simeq2.04\%,\quad P_{\rm out,1}\simeq9.24\%,\quad P_{\rm out,2}\simeq29.34\%.}
\end{equation}
These probabilities grow with increasing energy toward the continuum threshold in this example. A finite well has a \emph{finite} bound spectrum, in contrast to the infinite well. The plotted eigenfunctions have respectively zero, one, and two interior nodes.

\textbf{Step 11: Infinite-well limit.} With fixed $a,m$ and $U\to\infty$, the low-lying boundary amplitudes and penetration lengths tend to zero, while the energies measured from the bottom of the well approach the infinite-well values for width $L=2a$:
\begin{equation}
E_n+U\longrightarrow\frac{n^2\pi^2\hbar^2}{8ma^2},\qquad n=1,2,\ldots\quad\text{(fixed $n$)}.
\end{equation}
This reproduces the Dirichlet spectrum of Problem B5 after the appropriate shift in the energy zero.

\section{Consolidated answer sheet}
\begin{longtable}{@{}p{.11\linewidth}p{.81\linewidth}@{}}\toprule
Question & Main result \\\midrule\endhead
A1 & $\langle f_1,f_1\rangle=2/3$, $\langle f_2,f_2\rangle=2/5$, $\langle f_1,f_2\rangle=0$; $0\leq4/15$.\\
A2 & $[\hat p,\hat H]=0$, $[\hat x,\hat H]=\iu\hbar\hat p/m$ for a free particle.\\
A3 & $P_{n=1}=1/2+1/\pi\simeq81.83\%$, $P_{n=2}=1/2$.\\
A4 & $A=1/\sqrt2$, $\langle H\rangle=5\pi^2\hbar^2/(4mL^2)$, $\sigma_H=3\pi^2\hbar^2/(4mL^2)$.\\
A5 & In the verified \emph{second-edition} Griffiths variant, final probability after $A\!\to\!B\!\to\!A$ is $337/625$. Other editions have different Problem 3.27 statements (appendix).\\
B1 & $\Delta A\Delta B\geq\frac12\abs{\langle[A,B]\rangle}$; strengthened by the covariance term; $\Delta x\Delta p\geq\hbar/2$.\\
B2 & $R=(k_1-k_2)^2/(k_1+k_2)^2$, $T=4k_1k_2/(k_1+k_2)^2$ for $E>V_0$; $R+T=1$.\\
B3 & $R=1$, $T=0$, $\psi_{II}\propto\ee^{-\kappa x}$, $\delta=1/\kappa$ for a semi-infinite step with $E<V_0$.\\
B4 & $T=\left[1+\frac{V_0^2}{4E(V_0-E)}\sinh^2(\kappa a)\right]^{-1}$; for $\kappa a\gg1$, $T\sim16E(V_0-E)\ee^{-2\kappa a}/V_0^2$.\\
B5 & $E_n=n^2\pi^2\hbar^2/(2mL^2)$, $\psi_n=\sqrt{2/L}\sin(n\pi x/L)$.\\
B6 & $\Delta x\Delta p=\frac\hbar2\sqrt{n^2\pi^2/3-2}>\hbar/2$ for every $n\geq1$.\\
B7 & $k\tan(ka)=\kappa$ (even), $-k\cot(ka)=\kappa$ (odd); normalized states have $N_n=(a+1/\kappa_n)^{-1/2}$; at least three states iff $P>\pi$.\\\bottomrule
\end{longtable}

\appendix
\section{Griffiths Problem 3.27 in other editions}\label{app:griffiths_variants}
The uploaded assignment does \emph{not} specify which edition of \emph{Introduction to Quantum Mechanics} is intended. To prevent a misleading one-to-one identification, this appendix solves the two other verified Problem 3.27 versions. The results are alternative edition-specific reference solutions, \emph{not} additional assigned numbered questions. Problem A5 in the main text supplies the full 2005 version.

\subsection{First edition (1995): Gaussian-weighted polynomial space and Gram--Schmidt}
The 1995 Problem 3.27 concerns the finite-dimensional space of Gaussian-weighted polynomials
\begin{equation}
\Hil_N=\left\{p(x)\ee^{-x^2/2}:\deg p<N\right\},\qquad x\in\R,
\end{equation}
and asks for verification of its inner-product structure and Gram--Schmidt orthonormalization of the first four natural monomial-Gaussian functions \cite{griffiths1}.

\textbf{Step 1: Vector space and inner product.} For $f=p\ee^{-x^2/2}$, $g=q\ee^{-x^2/2}$ and scalars $\alpha,\beta\in\mathbb C$,
\begin{equation}
\alpha f+\beta g=(\alpha p+\beta q)\ee^{-x^2/2}\in\Hil_N.
\end{equation}
The ordinary $L^2(\R)$ scalar product is finite since
\begin{equation}
\langle f,g\rangle=\int_{-\infty}^{\infty}p^*(x)q(x)\ee^{-x^2}\,\dd x,
\end{equation}
and every polynomial multiplied by $\ee^{-x^2}$ is integrable. Conjugate symmetry, complex linearity in the second factor, and positivity follow from the $L^2$ inner product. A zero norm forces $p=0$ almost everywhere, hence as a polynomial $p\equiv0$. Thus $\Hil_N$ is an $N$-dimensional inner-product space and is automatically complete.

\textbf{Step 2: Gaussian moments.} Let $g_j(x)=x^j\ee^{-x^2/2}$ for $j=0,1,2,3$. The needed even moments are
\begin{equation}
\int_\R\ee^{-x^2}\,\dd x=\sqrt\pi,\quad
\int_\R x^2\ee^{-x^2}\,\dd x=\frac{\sqrt\pi}{2},\quad
\int_\R x^4\ee^{-x^2}\,\dd x=\frac{3\sqrt\pi}{4},\quad
\int_\R x^6\ee^{-x^2}\,\dd x=\frac{15\sqrt\pi}{8}.
\end{equation}
All odd Gaussian moments vanish by parity. The formulas follow by differentiating $I(\alpha)=\int_\R\ee^{-\alpha x^2}\dd x=\sqrt{\pi/\alpha}$ with respect to $\alpha$.

\textbf{Step 3: First two orthonormal functions.} $\|g_0\|^2=\sqrt\pi$ gives
\begin{equation}
\phi_0(x)=\pi^{-1/4}\ee^{-x^2/2}.
\end{equation}
By odd parity, $\langle g_0,g_1\rangle=0$. Since $\|g_1\|^2=\sqrt\pi/2$,
\begin{equation}
\phi_1(x)=\sqrt2\,\pi^{-1/4}x\ee^{-x^2/2}.
\end{equation}

\textbf{Step 4: Third function.} The projection of $g_2$ along $g_0$ is $(\langle g_0,g_2\rangle/\langle g_0,g_0\rangle)g_0=\frac12g_0$. Thus the orthogonalized vector is
\begin{equation}
u_2(x)=\left(x^2-\frac12\right)\ee^{-x^2/2}.
\end{equation}
Its squared norm is $\int_\R(x^2-\frac12)^2\ee^{-x^2}\dd x=\sqrt\pi/2$ by the moments above, hence
\begin{equation}
\phi_2(x)=\frac{2x^2-1}{\sqrt2\,\pi^{1/4}}\ee^{-x^2/2}.
\end{equation}

\textbf{Step 5: Fourth function.} Since $g_3$ is odd, its projections onto the even $\phi_0,\phi_2$ vanish. The $g_1$ projection coefficient is
\begin{equation}
\frac{\langle g_1,g_3\rangle}{\langle g_1,g_1\rangle}
=\frac{3\sqrt\pi/4}{\sqrt\pi/2}=\frac32.
\end{equation}
Therefore
\begin{equation}
u_3(x)=\left(x^3-\frac32x\right)\ee^{-x^2/2}.
\end{equation}
Its squared norm is
\begin{equation}
\|u_3\|^2=\int_\R(x^6-3x^4+\tfrac94x^2)\ee^{-x^2}\,\dd x
=\frac{15\sqrt\pi}{8}-\frac{9\sqrt\pi}{4}+\frac{9\sqrt\pi}{8}=\frac{3\sqrt\pi}{4}.
\end{equation}
Thus
\begin{equation}
\phi_3(x)=\frac{2x^3-3x}{\sqrt3\,\pi^{1/4}}\ee^{-x^2/2}.
\end{equation}
\begin{resultbox}
\textbf{First-edition answer.}
\begin{align*}
\phi_0(x)&=\pi^{-1/4}\ee^{-x^2/2},\\
\phi_1(x)&=\sqrt2\,\pi^{-1/4}x\ee^{-x^2/2},\\
\phi_2(x)&=\frac{2x^2-1}{\sqrt2\,\pi^{1/4}}\ee^{-x^2/2},\\
\phi_3(x)&=\frac{2x^3-3x}{\sqrt3\,\pi^{1/4}}\ee^{-x^2/2}.
\end{align*}
\end{resultbox}
They satisfy $\langle\phi_j,\phi_k\rangle=\delta_{jk}$. These are the first four normalized Hermite functions: in general $\phi_n(x)=(\pi^{1/4}\sqrt{2^nn!})^{-1}H_n(x)\ee^{-x^2/2}$. Their appearance connects Gram--Schmidt orthogonalization with the quantum harmonic oscillator.

\subsection{Third edition (2018): spectral decomposition and exponentials}
The third-edition Problem 3.27 asks to express an operator in its orthonormal eigenbasis and show that a spectral definition of an operator function agrees with the operator power series for the exponential \cite{griffiths3}.

\textbf{Step 1: Complete orthonormal eigenbasis.} Suppose $\hat Q\ket{e_n}=q_n\ket{e_n}$ and $\langle e_m|e_n\rangle=\delta_{mn}$, with completeness $\sum_n\ket{e_n}\bra{e_n}=\id$. For an arbitrary state $\ket\alpha$ in the domain,
\begin{align}
\hat Q\ket\alpha
&=\hat Q\sum_n\ket{e_n}\braket{e_n}{\alpha}
=\sum_nq_n\ket{e_n}\braket{e_n}{\alpha}\\
&=\left(\sum_nq_n\ket{e_n}\bra{e_n}\right)\ket\alpha.
\end{align}
Since the equality holds on the domain for every $\ket\alpha$,
\begin{equation}\boxed{\hat Q=\sum_nq_n\ket{e_n}\bra{e_n}.}\end{equation}

\textbf{Step 2: Operator exponential.} Since $\hat Q^j\ket{e_n}=q_n^j\ket{e_n}$ for all nonnegative integers $j$, the power series gives, first for finite-dimensional or bounded operators,
\begin{align}
\ee^{\hat Q}&=\sum_{j=0}^{\infty}\frac{\hat Q^j}{j!}
=\sum_{j=0}^{\infty}\frac1{j!}\sum_nq_n^j\ket{e_n}\bra{e_n}\\
&=\sum_n\left(\sum_{j=0}^{\infty}\frac{q_n^j}{j!}\right)\ket{e_n}\bra{e_n}
=\boxed{\sum_n\ee^{q_n}\ket{e_n}\bra{e_n}}.
\end{align}
For unbounded operators, this statement requires the appropriate spectral-functional-calculus definition and domain; formal interchange of all unbounded series is not always justified. A central physical example is the unitary time-evolution operator of a time-independent Hamiltonian,
\begin{equation}
\boxed{\hat U(t)=\ee^{-\iu\hat Ht/\hbar}=\sum_n\ee^{-\iu E_nt/\hbar}\ket n\bra n}
\end{equation}
in a complete discrete energy eigenbasis. This shows why each stationary state carries the phase $\ee^{-\iu E_nt/\hbar}$.

\section{Numerical and figure verification}\label{app:numerical}
The figures in B7 were generated with the companion Python file \texttt{make\_figures.py}, which uses only a one-dimensional root finder, exact piecewise eigenfunctions, and vector graphics. The numerical solver evaluates
\begin{equation}
F_n(z)=z+\arcsin(z/P)-\frac{(n+1)\pi}{2}=0,\qquad 0<z<P,
\end{equation}
with $P=4$. The unique roots $z_n$ give $q_n=\sqrt{16-z_n^2}$ and dimensionless normalization constants $\mathcal N_n=(1+1/q_n)^{-1/2}$. Substitution into $F_n$ gives residuals at floating-point accuracy; the exact normalization integral is one for every state by the analytic derivation. Quadrature over the plotted finite interval approximates, but does not exactly equal, the full-line normalization because it omits very small tails beyond the plot window. These checks help distinguish truncation of a figure from a failure of quantum-state normalization.

\section{Recommended bibliography and source notes}
\begin{thebibliography}{9}
\bibitem{assignment} Department of Physics \& Nanotechnology, SRM Institute of Science and Technology, \emph{Quantum Mechanics--I: Assignment II}, M.Sc. Physics / M.Sc. Applied Physics, Semester I, PPY25103J / PAP25103J (user-supplied course handout).
\bibitem{griffiths1} D. J. Griffiths, \emph{Introduction to Quantum Mechanics}, 1st ed. (Prentice Hall, 1995), Chapter 3, Problem 3.27. Edition-specific problem identification available at \url{https://faculty.washington.edu/seattle/physics441/online/Griffiths.pdf}.
\bibitem{griffiths2} D. J. Griffiths, \emph{Introduction to Quantum Mechanics}, 2nd ed. (Pearson, 2005), Chapter 3, Problem 3.27, ``Sequential measurements.'' Edition-specific problem identification at \url{https://iqti.iisc.ac.in/wp-content/uploads/2021/06/QM_Griffiths.pdf}.
\bibitem{griffiths3} D. J. Griffiths and D. F. Schroeter, \emph{Introduction to Quantum Mechanics}, 3rd ed. (Cambridge University Press, 2018), Chapter 3, Problem 3.27, spectral decomposition of operators; see \url{https://studylib.net/doc/28021667/griffiths---introduction-to-quantum-mechanics-3rd-ed-2018}.
\bibitem{shankar} R. Shankar, \emph{Principles of Quantum Mechanics}, 2nd ed. (Springer, 1994).
\bibitem{sakurai} J. J. Sakurai and J. Napolitano, \emph{Modern Quantum Mechanics}, 2nd ed. (Pearson, 2011).
\bibitem{reed_simon} M. Reed and B. Simon, \emph{Methods of Modern Mathematical Physics I: Functional Analysis}, rev. ed. (Academic Press, 1980).
\end{thebibliography}
\end{document}